Une randonneuse crie face à une falaise par une journée où la température de l'air est de \qty{\v{T}}{\degreeCelsius}. Elle entend l'écho \qty{\v{t}}{s} plus tard. La vitesse du son dans l'air est donnée par $v = \qty{331}{m/s}\sqrt{T/\qty{273}{K}}$, où $T$ est la température absolue.

\begin{enumerate}[label=\alph*)]
    \item Déterminez la vitesse du son dans l'air ce jour-là.
    \item À quelle distance de la falaise se trouve la randonneuse ?
    \item Au même endroit, par une journée où il fait \qty{\v{T2}}{\degreeCelsius}, après combien de temps entendrait-elle l'écho ?
\end{enumerate}
